\documentclass[10pt,conference]{IEEEtran}
\usepackage{cite}
\usepackage{amsmath,amssymb}
\usepackage{booktabs}
\usepackage{array}
\usepackage{url}
\usepackage[T1]{fontenc}
\renewcommand{\arraystretch}{1.08}
\begin{document}

\title{Where Did the Big Bang Start?\\
A Design-Only Inference Framework for the Origin of Spacetime}

\author{Survey--Idea--ExperimentDesign--Author Research Workflow}

\maketitle

\begin{abstract}
The everyday wording of ``Where did the Big Bang start?'' suggests an explosion at a position in an already-existing space. This paper separates that intuition from the testable content of modern cosmology. In a homogeneous and isotropic Friedmann--Lema\^itre--Robertson--Walker (FLRW) spacetime, cosmic expansion changes distances between comoving locations across the spatial slice; it does not select a spatial center. Observations supporting a hot, dense early universe and a model-inferred age near 13.8 Gyr therefore do not identify a location of the Big Bang or directly observe an initial instant. Inflationary descriptions, classical past-incompleteness theorems, bounce scenarios, and quantum-initial-state proposals address different parts of the extrapolation problem and must not be conflated. We propose the ORIGIN-OF-SPACETIME-INFERENCE-BENCHMARK, a preregistration-ready, design-only protocol for comparing only those early-universe scenario classes that expose explicit background histories, perturbation initial states, and forward predictions for cosmic microwave background (CMB), large-scale-structure, expansion-history, curvature, topology, or primordial-tensor observables. The protocol uses frozen parameter ledgers, held-out channels, synthetic injections, posterior-predictive checks, and foreground/prior robustness analyses. Its strongest direct conclusion is conditional but clear: under the empirically successful homogeneous FLRW description, the Big Bang was not located at one place in the universe. The physical completion of the earliest regime remains unresolved and requires discriminating observable predictions rather than a verbal account of ``before time.'' No new data analysis, posterior, model ranking, or empirical result is claimed here.
\end{abstract}

\begin{IEEEkeywords}
Big Bang cosmology, FLRW spacetime, cosmic microwave background, inflation, past incompleteness, bouncing cosmology, Bayesian inference, reproducible research.
\end{IEEEkeywords}

\section{Introduction}
\label{sec:introduction}

The question ``Where did the Big Bang start?'' is scientifically valuable only after its spatial presupposition is made explicit. A firecracker explodes at a point of a pre-existing room, so every later fragment can in principle be traced back to that point. The standard cosmological model is not this picture. It models the large-scale universe as a dynamical spacetime. The successful hot Big Bang account connects cosmic expansion, primordial nucleosynthesis, and the CMB to a hot, dense early epoch; it does not supply a camera image of a first event or a coordinate in an external space.

This distinction matters because several separate questions are routinely compressed into one sentence. Is there a preferred location on a spatial slice? How does the scale factor evolve? Does a classical spacetime admit arbitrarily extended past-directed geodesics? Which physical mechanism prepares the perturbations seen in the CMB? And does a mathematical boundary justify a statement about an absolute beginning? These questions have different assumptions, data access, and failure modes. Treating any one of them as an answer to all the others produces stronger claims than the evidence warrants.

Planck's analysis of CMB anisotropies constrains a six-parameter flat Lambda cold-dark-matter baseline and yields a model-dependent cosmic age close to 13.8 Gyr~\cite{planck2020}. That result is powerful evidence for a particular late-time-to-recombination cosmological description, but it is neither an observation of $t=0$ nor a measurement of a spatial starting point. Guth's inflationary proposal addresses the horizon and flatness problems and supplies a framework for primordial perturbations~\cite{guth1981}; it is not by itself a verified account of why the entire spacetime exists. Likewise, the past-incompleteness result of Borde, Guth, and Vilenkin identifies the need for physics beyond a sufficiently expanding or inflating description at a past boundary~\cite{bgv2003}. It does not prove creation from nothing, locate a center, or uniquely determine a preceding state.

The purpose of this report is therefore twofold. First, it gives a technically precise answer to the location question within the standard homogeneous model. Second, it turns the remaining origin question into a tractable model-comparison research program. The proposed benchmark does not reward a scenario for telling the most compelling narrative. It requires a declared forward map from early-universe variables to observables and then evaluates whether a prediction survives out-of-sample, foreground, and prior-sensitivity tests.

\section{Survey: What Observations and Theory Establish}
\label{sec:survey}

\subsection{Empirical scope of the hot Big Bang}

The hot Big Bang is best understood as an empirically constrained expansion history, not as a literal detonation. CMB observations provide a snapshot of the plasma after photons decoupled from matter, while galaxy-distance measurements, baryon acoustic oscillations, and abundance measurements constrain complementary portions of the cosmic history. In the baseline interpretation, these observations are mutually consistent with an early hot and dense state followed by expansion and structure formation. They justify extrapolation over an immense but finite physical interval only within the tested model family and its stated assumptions.

The inferred age is a derived parameter. It follows after a cosmological parameterization, a likelihood, calibration information, and a relation between model time and observables are adopted. It must not be described as the elapsed time since a directly observed explosion point. The difference is not semantic: a fitted parameter can change when a model, prior, data combination, or systematic treatment changes, whereas a location in an external room would be a different kind of observable entirely.

\subsection{FLRW geometry and the absence of a center}

On sufficiently large scales, the conventional background is an FLRW line element,
\begin{equation}
ds^2=-c^2dt^2+a^2(t)\left[\frac{d\chi^2}{1-k\chi^2}+\chi^2\left(d\theta^2+\sin^2\theta\,d\phi^2\right)\right],
\label{eq:flrw}
\end{equation}
where $t$ is cosmological proper time for comoving observers, $c$ is the speed of light, $a(t)$ is the scale factor, $(\chi,\theta,\phi)$ are comoving coordinates, and $k$ labels spatial-curvature sign and normalization. Equation~\eqref{eq:flrw} is a symmetry-based model of a spatially homogeneous and isotropic universe. Homogeneity means that no comoving point is physically distinguished by the background solution, while isotropy means that no direction is distinguished at a comoving point.

For neighboring comoving galaxies separated by a fixed coordinate interval $\Delta\chi$, their proper separation is
\begin{equation}
D(t)=a(t)\Delta\chi.
\label{eq:proper-distance}
\end{equation}
Thus the expansion statement is about the time dependence of $D(t)$ in \eqref{eq:proper-distance}. It does not require an external embedding space or a central source. Differentiating this relation gives the local recession law $\dot{D}=H D$ when the comoving separation is fixed, with the Hubble rate defined by
\begin{equation}
H(t)=\frac{\dot{a}(t)}{a(t)}.
\label{eq:hubble-rate}
\end{equation}
The dot in \eqref{eq:hubble-rate} denotes a derivative with respect to $t$. Every ideal comoving observer in the homogeneous background can describe other sufficiently nearby comoving observers as receding according to the same form. This reciprocity is the reason a local Hubble pattern cannot by itself identify a central observer.

An observational horizon is also not a center. It bounds the region from which light has had time to reach a particular observer. A second observer has a different horizon centered on that observer. The finite observable universe is therefore an observer-dependent causal region inside a model, not a map of the total universe with Earth at a privileged physical position.

\subsection{Early-universe mechanisms and classical boundaries}

The background expansion is commonly summarized by a Friedmann equation such as
\begin{equation}
H^2=\frac{8\pi G}{3}\rho-\frac{k c^2}{a^2}+\frac{\Lambda c^2}{3},
\label{eq:friedmann}
\end{equation}
where $G$ is Newton's constant, $\rho$ is total energy density, and $\Lambda$ is the cosmological constant. Equation~\eqref{eq:friedmann} has predictive content within general relativity and a specified stress-energy model. Its formal backward extrapolation to very high density is not proof that the mathematical singularity is an observed physical object. Classical singularity theorems instead establish geodesic incompleteness under particular causal, energy, and global assumptions~\cite{hawkingpenrose1970}.

Inflation adds a phase of accelerated expansion to address selected initial-condition problems and to generate a calculable primordial fluctuation spectrum~\cite{guth1981}. Bouncing and ekpyrotic constructions attempt a nonsingular or matched prior phase, but they introduce their own consistency, stability, matching, and observational challenges~\cite{brandenberger2017}. Quantum-initial-state or quantum-cosmology proposals may be scientifically comparable only when they provide a quantitative background and perturbation prescription. A proposal lacking a forward prediction is an interesting theoretical possibility but is underdetermined as an observational hypothesis.

\section{Idea: An Origin-of-Spacetime Inference Benchmark}
\label{sec:idea}

\subsection{Scientific reframing and objective}

The ORIGIN-OF-SPACETIME-INFERENCE-BENCHMARK replaces the unobservable request for a location before space with a sequence of testable questions: Which background and perturbation histories reproduce the measured relics of the early universe? Do alternatives make predictions that the reference model fails to make? Are apparent improvements robust to held-out data, nuisance modeling, and prior choices? The benchmark separates a model's observable utility from any stronger metaphysical interpretation of its extrapolated boundary.

The reference class $S0$ is a hot Big Bang late-time baseline with a phenomenological primordial spectrum. Candidate class $S1$ contains parameterized slow-roll or effective inflationary histories. Class $S2$ contains parameterized bounce or ekpyrotic histories with a declared matching surface. Class $S3$ contains quantum-initial-state models only when they expose explicit background and perturbation predictions. Class $S4$ contains curvature or topology extensions and is deliberately kept distinct from an origin-mechanism claim. This separation prevents a finite topology, nonzero curvature, or finite horizon from being mistaken for evidence of a central explosion.

\begin{table*}[!t]
\caption{Scenario classes and admissibility requirements in the proposed benchmark. ``Design only'' means that this report does not execute the comparison or claim outcomes.}
\label{tab:scenario-ledger}
\centering
\footnotesize
\begin{tabular}{p{0.08\textwidth}p{0.23\textwidth}p{0.31\textwidth}p{0.28\textwidth}}
\toprule
Class & Role & Required declared objects & Interpretation boundary \\
\midrule
$S0$ & Hot Big Bang baseline & Late-time parameters, nuisance model, phenomenological primordial spectrum, transfer calculation & Reference predictive model; not a claim about a first cause. \\
$S1$ & Inflationary scenarios & Background potential or effective history, reheating assumptions, scalar and tensor initial conditions & Can address horizon/flatness and perturbations; does not alone establish an ultimate beginning. \\
$S2$ & Bounce/ekpyrotic scenarios & Contracting/bounce history, matching prescription, stability and validity range, perturbation transfer & A candidate completion only if it has a stable forward prediction distinct from $S0$. \\
$S3$ & Quantum-initial-state scenarios & Boundary/state prescription, calculable background, perturbation spectrum, parameter priors & A verbal no-boundary or tunneling label alone is not numerically comparable. \\
$S4$ & Curvature/topology extensions & Curvature/topology parameters and predicted map/spectrum signatures & Geometry is tested separately and never treated as a coordinate center. \\
\bottomrule
\end{tabular}
\end{table*}

\subsection{Falsifiable hypotheses}

The benchmark has five operational hypotheses. H0 states that $S0$ accounts for the selected observations as well as a more complex scenario after complexity and prior sensitivity are included. H1 states that an alternative yields a posterior-predictive improvement on predeclared held-out observables. H2 states that a putative primordial feature weakens after foreground, look-elsewhere, and prior-volume controls. H3 states that a horizon, curvature estimate, or topology signature does not identify a Big Bang center in a homogeneous model. H4 states that past incompleteness requires a boundary description for a model class but does not select a unique physical origin. H5 states that statistical model comparison can compare predictive scenarios but cannot answer an undefined question about an observed time ``before time.''

These hypotheses are intentionally asymmetric. A null result does not prove that no pre-hot-Big-Bang physics exists; it limits the discriminating power of the tested data and parameterization. Conversely, a fit improvement does not establish an origin story unless it is replicated in held-out observables, survives the declared robustness checks, and is not reproducible by a comparably flexible baseline.

\section{Formal Cosmological Framework}
\label{sec:framework}

\subsection{Observable mapping}

The redshift assigned to a photon emitted at scale factor $a_{\mathrm{em}}$ and observed at $a_0$ follows
\begin{equation}
1+z=\frac{a_0}{a_{\mathrm{em}}},
\label{eq:redshift}
\end{equation}
where $z$ is redshift and the present scale factor is often normalized to $a_0=1$. Equation~\eqref{eq:redshift} relates measured wavelength shifts to the expansion model. It does not select a center because every comoving observer uses the same local relation under the FLRW assumptions.

For a specified expansion history, a comoving particle-horizon distance may be written as
\begin{equation}
\chi_{\mathrm{ph}}(t)=\int_{t_{\mathrm{init}}}^{t}\frac{c\,dt'}{a(t')}.
\label{eq:particle-horizon}
\end{equation}
In \eqref{eq:particle-horizon}, the lower limit $t_{\mathrm{init}}$ is a model boundary, not a directly observed first moment. A finite value constrains causal contact in that model. It cannot be used to infer a point in an external space, and divergent or finite versions of the integral depend on the assumed extension of $a(t)$.

The benchmark takes a scenario to be a tuple $M_i=\{B_i,I_i,F_i,\eta_i,\Pi_i\}$: background history $B_i$, perturbation initial-state prescription $I_i$, forward model $F_i$, nuisance/late-time variables $\eta_i$, and prior specification $\Pi_i$. The forward model propagates primordial scalar and tensor quantities through a transfer calculation and instrument/foreground response to the selected CMB and late-time observables. This representation makes it possible to identify which part of a conclusion comes from data, which from transfer physics, and which from a chosen initial-state ansatz.

\subsection{Inference quantities}

For data $D$, model parameters $\theta_i$, and scenario $M_i$, the posterior is
\begin{equation}
p(\theta_i\mid D,M_i)=\frac{p(D\mid\theta_i,M_i)\,p(\theta_i\mid M_i)}{p(D\mid M_i)}.
\label{eq:posterior}
\end{equation}
Here $p(D\mid\theta_i,M_i)$ is the likelihood, $p(\theta_i\mid M_i)$ is the declared prior, and the denominator is the model evidence. The latter is the prior-weighted likelihood integral,
\begin{equation}
p(D\mid M_i)=\int p(D\mid\theta_i,M_i)\,p(\theta_i\mid M_i)\,d\theta_i.
\label{eq:evidence}
\end{equation}
Equations~\eqref{eq:posterior} and \eqref{eq:evidence} clarify why a bare evidence or Bayes-factor number cannot be treated as an origin verdict: it changes with the parameter space, prior, likelihood, and model definition being compared.

For a held-out datum $D_{\mathrm{hold}}$, the predictive distribution is
\begin{align}
p(D_{\mathrm{hold}}\mid D_{\mathrm{train}},M_i)
&=\int p(D_{\mathrm{hold}}\mid\theta_i,M_i) \nonumber\\
&\quad\times p(\theta_i\mid D_{\mathrm{train}},M_i)\,d\theta_i.
\label{eq:posterior-predictive}
\end{align}
Equation~\eqref{eq:posterior-predictive} is the primary protection against a scenario that merely accommodates a tuned feature in the training channel. It evaluates whether the scenario predicts an independent observable after its parameters have been constrained without access to that observable.

\section{ExperimentDesign: A Preregistered Comparison Protocol}
\label{sec:experiment-design}

\subsection{Design status and research unit}

This section specifies a design-only experiment. It neither downloads a particular data release nor executes a likelihood evaluation. Consequently, no numerical preference, posterior interval, Bayes factor, anomaly significance, or physical conclusion beyond the symmetry reasoning in Section~\ref{sec:survey} is reported. The research unit is a fully declared scenario instance rather than a named cosmological slogan. An instance is eligible only when its background evolution, perturbation initial conditions, late-time connection, and observable likelihood are available at a stated numerical precision.

The intended data domains are CMB temperature and polarization angular spectra, CMB lensing, baryon acoustic oscillations, distance measures, large-scale structure, and primordial-abundance consistency. The exact release, sky masks, calibration treatment, beam model, covariance, foreground parameterization, and likelihood software must be versioned before fitting. This rule prevents a model comparison from silently changing as the infrastructure changes. It also distinguishes a theoretical prediction from a data-product convention.

\subsection{Frozen scenario and parameter ledger}

Before any production analysis, every scenario must submit a machine-readable ledger. The ledger records background variables; scalar and tensor spectrum parameters; reheating, sound-speed, matching, or transition parameters; late-time cosmological variables; nuisance variables; prior distributions; parameter transformations; theoretical validity domains; physicality checks; and computational settings. The baseline ledger is frozen alongside the alternative ledger. A new parameter cannot be introduced after a residual is inspected unless the study is explicitly restarted and the change is disclosed.

\begin{table}[!t]
\caption{Minimum preregistered ledger for a scenario instance. All entries are frozen before held-out evaluation.}
\label{tab:parameter-ledger}
\centering
\footnotesize
\setlength{\tabcolsep}{3pt}
\begin{tabular}{p{0.25\columnwidth}p{0.59\columnwidth}}
\toprule
Ledger item & Required content and reason \\
\midrule
Background & $a(t)$ or an equivalent solved history, field/content assumptions, and validity range; prevents an unexamined extrapolation. \\
Initial state & Scalar and tensor initial conditions with normalization and matching convention; makes the primordial prediction reproducible. \\
Transition physics & Reheating, bounce matching, sound speed, or effective-theory cutoff; exposes quantities that can mimic a feature. \\
Late-time and nuisance terms & Shared cosmological parameters, calibration, beams, masks, and foregrounds; enables matched treatment. \\
Prior file & Family, bounds, transformations, and motivation; makes evidence sensitivity auditable. \\
Forward implementation & Transfer-code version, numerical tolerance, solver failures, and likelihood wrapper; enables numerical replication. \\
Prediction list & Which temperature, polarization, lensing, spectral, curvature, topology, or tensor signatures are expected; prevents post hoc target selection. \\
\bottomrule
\end{tabular}
\end{table}

The physicality policy is as important as the sampler. A point that violates a declared stability condition, a positivity condition, an effective-theory range, or the model's matching assumptions is invalid rather than being assigned a convenient likelihood. Invalid points, failed solver calls, discarded chains, and nonconverged runs are retained in the run record. Otherwise, a numerically favorable but physically undefined region can create a false appearance of support.

\subsection{Data partitioning and held-out tests}

The study partitions observables by role. Development data are used only to validate the transfer code on synthetic skies, set fixed numerical tolerances, and test the ingestion pipeline. Training data constrain the declared parameter ledger. At least one predeclared channel remains held out. Suitable choices include polarization or lensing withheld from a temperature-led fit; a detector or frequency split withheld from feature selection; a multipole interval withheld from a search for oscillations or a cutoff; or low-redshift distance data withheld from an early-universe model choice. The partition must respect correlations: if a data product is correlated with a training product, a joint covariance or an explicitly independent split is required.

The validation set is not a source of exploratory tuning. A candidate feature selected after examining a particular low-multipole range, for example, must be treated as exploratory and evaluated on independently simulated or newly released data. This distinction directly addresses the look-elsewhere problem, in which many possible feature positions, widths, and shapes are searched but only the most striking fluctuation is reported.

\subsection{Forward calculation}

For each valid parameter vector, the workflow solves the background and perturbation equations, converts primordial quantities into transfer functions, predicts angular or spatial observables, applies the detector and foreground model, and evaluates the likelihood. A reference transfer implementation is checked against an independent implementation or independently validated mode wherever feasible. Numerical agreements are assessed on a grid that includes baseline-like, boundary, and feature-bearing points, not merely at a single favorable parameter vector.

The same masks, likelihood family, calibration priors, covariance treatment, and late-time parameterization are used for the baseline and each physically compatible alternative. If a scenario genuinely requires a different likelihood or transfer structure, the departure is declared and a matched sensitivity analysis is performed. Differences in analysis convenience must never be allowed to masquerade as physical predictive gain.

\section{Scenario Tests, Systematics, and Analysis}
\label{sec:analysis}

\subsection{No-center geometry simulation}

The first task is intentionally simple and conceptually decisive. Generate synthetic source catalogs in flat, open, and closed FLRW backgrounds. Place virtual comoving observers at multiple coordinates, propagate their past light cones, and estimate local recession, redshift-distance, horizon, and angular-diameter relations. The expected invariance is not that all catalog realizations are identical; finite sampling, curvature, and observational selection create differences. The expected result is that after coordinate changes and controlled selection effects, the homogeneous background supplies no location that is privileged as ``the Big Bang center.''

This simulation is a test of the interpretation of the assumed symmetry model, not a new measurement that chooses among all cosmologies. A distinct central-origin claim can be investigated only by writing down a specified inhomogeneous spacetime, providing its forward predictions, and comparing it with data under the same standards. A coordinate origin used in numerical integration is never entered into the likelihood as an empirical variable.

\subsection{CMB and primordial-spectrum tests}

The primary likelihood tasks compare predicted CMB temperature, polarization, and lensing observables. Candidate scenarios may differ in spectral tilt, running, oscillatory structure, cutoff behavior, tensor amplitudes, or non-Gaussianity. Each search domain must be specified before inspecting the data. For example, a feature model declares the permitted frequency/width interval and its prior; a tensor model declares its relation to the scalar sector; and a bounce model declares the transfer across its matching surface. A flexible phenomenological feature is not automatically evidence for the underlying mechanism that inspired it, because several mechanisms can reproduce similar spectra.

The analysis reports a vector rather than a single origin score: likelihood behavior, convergence diagnostics, predictive residuals, posterior-predictive discrepancy, evidence or a stated approximation, complexity, parameter degeneracy, prior sensitivity, predictive coverage, and resource cost. A maximum-likelihood improvement alone is insufficient, because a more flexible model can fit fluctuations without predicting them. A scenario earns observational utility only when it improves a predeclared held-out quantity and that improvement persists under the robustness matrix described below.

\subsection{Expansion, curvature, and topology tests}

Expansion-history consistency uses CMB-derived constraints together with baryon acoustic oscillations, distance indicators, large-scale structure where appropriate, and primordial-nucleosynthesis consistency. These probes have different systematics and redshift leverage. Their combination is valuable precisely because a background history that fits one product can fail another, but correlations and selection functions must be propagated rather than ignored.

Curvature and topology are separate tests. Curvature modifies distance-redshift and mode relations; nontrivial topology can motivate specific correlation or pattern searches. Neither possibility changes the fact that a homogeneous FLRW model lacks a special spatial point. A finite closed spatial section need not have a center within that section, just as a two-dimensional sphere has no preferred location on its surface. The analogy is only topological intuition; quantitative conclusions require the explicit three-dimensional metric and data model.

\begin{table*}[!t]
\caption{Robustness matrix for the benchmark. A claimed scenario advantage must remain interpretable across the applicable rows.}
\label{tab:robustness-matrix}
\centering
\footnotesize
\begin{tabular}{p{0.18\textwidth}p{0.35\textwidth}p{0.39\textwidth}}
\toprule
Threat & Predeclared perturbation & Required diagnostic \\
\midrule
Foreground leakage & Vary foreground families, masks, frequency selections, and calibration nuisance priors. & Compare residual morphology and held-out coverage; report whether a feature tracks a foreground choice. \\
Prior-volume effect & Vary motivated prior families and widths without changing data. & Report evidence and posterior-predictive sensitivity; do not quote a prior-dependent ranking as absolute. \\
Feature selection & Withhold a multipole interval, detector split, or synthetic realization from selection. & Evaluate the frozen feature on the held-out channel and account for the declared search domain. \\
Transfer/numerical error & Use an independent implementation or validated mode, tolerance sweeps, and boundary tests. & Store deviations, solver failures, and convergence diagnostics; repeat only impacted calculations. \\
Late-time degeneracy & Change compatible late-time nuisance treatment and add independent background probes. & Report parameter shifts and whether an early feature is exchanged for late-time freedom. \\
Model misspecification & Inject known oscillations, cutoffs, curvature, and null skies into synthetic data. & Quantify recovery, false-positive behavior, and calibration of posterior summaries. \\
\bottomrule
\end{tabular}
\end{table*}

\subsection{Synthetic injection and calibration}

Synthetic injection studies are performed before interpreting the real-data comparison. Simulated skies and late-time summaries are generated from baseline-like and alternative-like parameter points, including null signals, controlled oscillations, cutoffs, curvature perturbations, and topology templates when defined. The analysis pipeline is then run unchanged. This exercise measures whether the sampler recovers a known signal, whether its uncertainties cover the generating value at the intended frequency, and whether foreground or numerical mistakes induce a false feature.

Simulation-based calibration does not prove that the universe was generated by the model. It tests an inferential pipeline under data generated from declared assumptions. When recovery fails, the relevant model family, likelihood, or numerical implementation must be repaired before real data are used for a mechanism claim. The synthetic records are released with seeds, injected parameter files, software versions, and success/failure labels.

\subsection{Decision rules}

The decision rules are deliberately limited. A scenario may be described as having improved observational utility only if its frozen pipeline shows posterior-predictive improvement on held-out observables, its result survives relevant rows of Table~\ref{tab:robustness-matrix}, and the comparison exposes its additional degrees of freedom. A scenario may be described as a candidate physical mechanism only if it supplies a discriminating signature that the baseline and competing alternatives cannot reproduce under matched nuisance treatment. Otherwise it is reported as inconclusive, underdetermined, or unsupported by the selected observables.

No decision rule maps data to an observed ``before time.'' Past-directed incompleteness and a finite domain of validity identify where a model stops describing the past; they do not turn the missing extension into a measured event. This boundary discipline is not a weakness of the benchmark. It is what allows a direct answer to a well-defined location question while retaining intellectual honesty about unresolved initial physics.

\section{Reproducibility, Governance, and Interpretation Safety}
\label{sec:reproducibility}

\subsection{Reproducible release package}

An executable benchmark release should contain scenario definitions, parameter ledgers, prior files, data-release identifiers, masks, likelihood wrappers, transfer-code versions, compilation settings, random seeds, sampler settings, chain diagnostics, synthetic injections, and an append-only event log. Public data are referenced by stable identifiers; restricted products are described by access conditions rather than redistributed. Posterior samples may be released only when licenses permit. Every figure and table produced by a future execution must be traceable to a run manifest that records exact input versions and an output checksum.

The release distinguishes four epistemic layers. Layer one contains observations and calibrated data products. Layer two contains parameter inferences conditional on a likelihood, calibration, and model. Layer three contains extrapolations into epochs not directly observed. Layer four contains theoretical completions outside the tested domain. The report and any downstream communication must label which layer supports each sentence. In particular, a data-supported parameter estimate cannot be silently promoted to a verified description of a primordial or pre-geometric regime.

\subsection{Computational governance}

Scenario authors must disclose computational budgets, failed attempts, filtering criteria, and any manual intervention. Sampler convergence is necessary but not sufficient: a converged chain can faithfully sample a misspecified likelihood. Conversely, a model with a computationally difficult but physically required region should be represented through improved algorithms or transparent approximation studies, not through deletion of unfavorable parameter space.

Data governance includes version pinning, blinded held-out channels, review of selection effects, and a preregistered procedure for unblinding. A compact governance rule is useful: no analyst may change a prior, feature range, mask, transfer tolerance, or model parameterization after seeing held-out performance without creating a new run identifier. The original run remains visible. This prevents a history of adaptive choices from being misrepresented as a single confirmatory test.

\section{Expected Observation Patterns and Limits}
\label{sec:expected-patterns}

Table~\ref{tab:expected-patterns} summarizes interpretation patterns that a future execution may encounter. It intentionally states outcomes conditionally. The table does not contain results from the present report.

\begin{table*}[!t]
\caption{Conditional interpretation patterns for a future benchmark execution. These are protocol outcomes, not empirical findings of this paper.}
\label{tab:expected-patterns}
\centering
\footnotesize
\begin{tabular}{p{0.27\textwidth}p{0.30\textwidth}p{0.36\textwidth}}
\toprule
Observed protocol pattern & Permitted inference & Inference that remains impermissible \\
\midrule
$S0$ predicts held-out channels as well as alternatives after robustness checks. & Selected data do not demonstrate added utility for the tested alternatives. & The baseline is the unique theory of the universe's origin. \\
An alternative improves held-out predictions under matched nuisance models and stable priors. & That parameterized scenario has comparative predictive utility for the declared observables. & The mechanism is the final or unique account of why spacetime exists. \\
Feature support disappears under frequency, mask, or prior variations. & The feature is not robust under the chosen protocol. & The entire early-universe framework is falsified. \\
An explicit inhomogeneous model outperforms FLRW on defined data. & The homogeneous baseline is limited for that defined comparison. & A coordinate point is thereby established as the Big Bang's location without further geometric analysis. \\
Past-incompleteness assumptions apply to a scenario. & That scenario needs a boundary description or additional past physics. & A material singularity, a spatial center, creation from nothing, or an observed pre-time event has been proved. \\
\bottomrule
\end{tabular}
\end{table*}

The limits in Table~\ref{tab:expected-patterns} follow from the mapping between theory and observable. CMB photons were released long after any hypothesized first instant. Even an exceptionally precise primordial-spectrum measurement would constrain a model's initial conditions through a transfer chain; it would not photograph a singularity. The farther a conclusion is extrapolated from the observable epoch, the more explicitly its theoretical assumptions must be stated.

The benchmark nevertheless advances the origin discussion. It prevents an untestable location claim from blocking research on testable remnants. A no-center conclusion becomes a geometrical implication that can be explained clearly. Candidate pre-hot-Big-Bang histories become scientific when they specify their perturbations and observational signatures. Past-boundary theorems remain important because they constrain the domain of a model, while not being misused as a full ontology of the universe.

\section{Threats to Validity and Falsifiability}
\label{sec:validity}

\subsection{Conceptual threats}

The first threat is a category mistake: taking an analogy that is useful for visualizing recession and treating it as a literal spacetime model. A balloon analogy may convey that every point on a two-dimensional surface recedes from every other point as the surface expands, but it also tempts a reader to identify the balloon's unphysical embedding-space center with a center of the universe. The benchmark avoids that transfer by working from the metric in \eqref{eq:flrw} and by declaring an external spatial embedding absent unless a competing model explicitly introduces one.

The second threat is equivocation between a model boundary and a physical event. General relativity has a well-defined domain of application, and classical singularity theorems establish conditions under which geodesics cannot be extended within a classical spacetime~\cite{hawkingpenrose1970}. The conclusion concerns the incompleteness of a description under its assumptions. It does not measure a physical point of infinite density, determine the relevant quantum-gravitational degrees of freedom, or prove that ordinary temporal language extends through the boundary. The benchmark requires each scenario to identify its validity range precisely so that an extrapolation is never silently converted into an observation.

The third threat is theory-label ambiguity. ``Inflation,'' ``bounce,'' ``no boundary,'' and ``emergent universe'' each name a family with variants rather than a single probability distribution. A comparison that uses a broad and flexible alternative but a narrow baseline can manufacture apparent support through unequal opportunity to fit. Conversely, a rigid toy implementation can incorrectly rule out an entire family. The remedy is to compare declared scenario instances, report what the instance does not cover, and avoid global statements about a family that was not fully parameterized.

\subsection{Statistical threats}

Cosmological data contain cosmic variance, instrumental noise, masking, foreground residuals, calibration uncertainty, and correlations across derived products. Low-multipole CMB quantities are particularly vulnerable to post-selection because there are relatively few modes and a great many imaginable patterns that may be noticed afterward. A visually unusual map feature is therefore not an origin signature without a predeclared statistic, null distribution, and independent test channel. The held-out design and synthetic null ensembles turn this principle into an executable constraint rather than an editorial preference.

Bayesian evidence also requires care. Equation~\eqref{eq:evidence} averages likelihood over a prior volume. That is appropriate when the prior is part of the declared model, but it means evidence values are conditional on that prior and cannot be exported as prior-free facts. The protocol reports sensitivity to justified prior alternatives and pairs any evidence calculation with posterior-predictive performance. Information criteria may be reported as complementary approximations, but their assumptions and treatment of effective complexity must be documented.

The problem of nonidentifiability is separate from statistical uncertainty. Two physically distinct histories can project to nearly the same scalar spectrum after transfer, and a primordial feature can be degenerate with calibration or foreground freedom. Narrow credible intervals in a restricted parameterization do not resolve this structural ambiguity. A credible research claim must state which degeneracies were tested, which independent observable could break them, and which conclusion stays conditional because an observable mapping is shared.

\subsection{Falsifiability and failure reporting}

A model is exposed to falsification when it restricts an observable distribution before the relevant data are examined. In this design, examples include a constrained relation among scalar tilt, tensor amplitude, non-Gaussianity, or a feature phase; a predicted consistency relation after reheating assumptions; or a topology pattern with a fixed template family. The prediction need not be simple, but it must be computable with declared uncertainties. A scenario that can absorb every residual by adding an unconstrained feature is not made stronger by its flexibility; it has avoided a test.

Failure reporting is correspondingly a first-class output. The release records rejected parameter points, unstable solutions, nonconverged sampling, failed synthetic recoveries, and results that disappear under an alternate foreground model. These outcomes identify where the theory or computation needs work. Suppressing them would bias the comparison toward narratives that happened to be numerically convenient. The success criterion is not confirmation of a cosmic beginning story; it is a transparent map from a proposed early-universe model to what would count against it.

\subsection{Scope of a future conclusion}

If a future execution finds no discriminating signature, the appropriate statement is that the selected observations and protocol do not distinguish the tested scenarios at the reported sensitivity. If it finds a robust signature, the appropriate statement is that a specified scenario class has greater predictive support for that observable set. Both statements are narrower than ``we know how the universe began.'' The distinction protects the value of the scientific result while acknowledging that a physical theory may be incomplete outside the presently observable domain.

\section{Discussion}
\label{sec:discussion}

The central answer can now be stated precisely. Within a standard homogeneous FLRW model, the Big Bang did not start at a distinguished place in the universe. The model represents an early state of the spatial universe in which separations between comoving locations were much smaller, and then grow according to the scale factor. Saying that ``everywhere was once hot and dense'' is closer to the model than saying that matter flew outward from a central point. This does not mean that every imaginable cosmology has no center; it means that a center is absent from the standard homogeneous background and would require a separately specified, testable departure.

The word ``start'' retains legitimate scientific uses, but they must be qualified. It can mean the earliest epoch to which a particular classical solution is extrapolated. It can mean a past boundary associated with geodesic incompleteness under stated conditions. It can mean the beginning of a hot, radiation-dominated phase in a parameterized scenario. These are not interchangeable. In particular, a statement that inflation does not by itself provide a complete past extension~\cite{bgv2003} does not identify what precedes it, and a bouncing model review~\cite{brandenberger2017} does not establish that a bounce occurred.

The recommended scientific posture is neither to declare the question meaningless nor to treat a popular visual metaphor as data. The location wording is a category error under the standard symmetry assumptions, whereas the causal and mechanism questions motivate research. The benchmark offers a way to investigate them with ordinary scientific disciplines: clear hypotheses, explicit model classes, shared nuisance treatment, synthetic calibration, held-out prediction, and reproducible artifacts.

Future extensions may incorporate new CMB polarization data, spectral distortions, gravitational-wave backgrounds, or improved structure surveys. Such extensions should enlarge observational access, not weaken the boundary discipline. Each new observable must enter through a frozen forward model, documented covariance, and a stated relation to the scenario's perturbation physics. A new data channel can discriminate among defined models; it cannot by itself answer a question whose terms have not been translated into an observable spacetime construction.

\section{Conclusion}
\label{sec:conclusion}

This four-stage research workflow reaches a direct but bounded conclusion. The observed hot, expanding universe supports a model-inferred early history, yet does not reveal a spatial point from which the universe exploded. In homogeneous FLRW cosmology, the Big Bang is not located at one place: expansion occurs between comoving locations throughout the spatial slice, and a horizon or coordinate origin is not a physical center. The mechanism responsible for the earliest applicable state remains open.

The ORIGIN-OF-SPACETIME-INFERENCE-BENCHMARK supplies a concrete next step. It compares only parameterized scenarios with explicit backgrounds, initial states, parameter ledgers, and forward predictions. Its held-out, prior-sensitive, and systematics-aware protocol can establish comparative predictive performance, expose underdetermined proposals, and delimit the reach of an extrapolation. It cannot observe a first instant, identify ``before time,'' or convert a past-boundary theorem into proof of a unique cosmic origin. This paper is therefore a reproducible research design, not an empirical model-selection result.

\appendices
\section{Minimal Scenario and Run Manifest}
\label{app:manifest}

For each submitted scenario, the minimum manifest fields are: scenario identifier and version; reference publications; background equations and validity regime; perturbation-state prescription; complete parameter ledger; prior file checksum; shared late-time and nuisance definitions; transfer/likelihood software versions; training and held-out data identifiers; blinding status; synthetic-injection identifiers; sampler configuration; numerical tolerances; convergence rule; physicality and invalid-point records; output artifact hashes; and an event log. The manifest includes a plain-language interpretation boundary: ``This run compares declared observable predictions and does not infer a spatial Big Bang center, a verified first instant, or a state before time.''

\section{Protocol Checklist}
\label{app:checklist}

Before fitting, verify that the scenario has an explicit forward map; freeze the scenario and prior ledgers; version data and masks; validate on synthetic data; and define held-out channels. During fitting, retain failed evaluations, use matched nuisance models, monitor convergence, and avoid retuning on held-out channels. Before interpretation, execute the robustness matrix, report parameter shifts and prior sensitivity, inspect posterior-predictive residuals, and classify a model with no discriminating forward prediction as underdetermined. Before release, archive manifests, code revisions, permitted samples, synthetic inputs, diagnostics, and an interpretation-layer statement.

\begin{thebibliography}{00}

\bibitem{planck2020}
Planck Collaboration, ``Planck 2018 results. VI. Cosmological parameters,'' \emph{Astronomy \& Astrophysics}, vol. 641, Art. no. A6, 2020, doi: 10.1051/0004-6361/201833910.

\bibitem{guth1981}
A. H. Guth, ``Inflationary universe: A possible solution to the horizon and flatness problems,'' \emph{Physical Review D}, vol. 23, no. 2, pp. 347--356, 1981, doi: 10.1103/PhysRevD.23.347.

\bibitem{hawkingpenrose1970}
S. W. Hawking and R. Penrose, ``The singularities of gravitational collapse and cosmology,'' \emph{Proceedings of the Royal Society A}, vol. 314, no. 1519, pp. 529--548, 1970, doi: 10.1098/rspa.1970.0021.

\bibitem{bgv2003}
A. Borde, A. H. Guth, and A. Vilenkin, ``Inflationary spacetimes are incomplete in past directions,'' \emph{Physical Review Letters}, vol. 90, Art. no. 151301, 2003, doi: 10.1103/PhysRevLett.90.151301.

\bibitem{brandenberger2017}
R. H. Brandenberger and P. Peter, ``Bouncing cosmologies: Progress and problems,'' \emph{Foundations of Physics}, vol. 47, pp. 797--850, 2017, doi: 10.1007/s10701-016-0057-0.

\bibitem{weinberg2008}
S. Weinberg, \emph{Cosmology}. Oxford, U.K.: Oxford Univ. Press, 2008.

\end{thebibliography}

\end{document}
